\label{chap:Results}
This chapter presents a structured evaluation of the seismic monitoring 
pipeline under multiple model configurations and parameter settings. The 
analysis is structured hierarchically, beginning with the picker module, 
followed by the associator and location modules, and concluding with the 
magnitude estimation and catalog recall assessments. We interpret the results 
in the broader context of operational seismology and ML based detection, 
highlighting the trade-offs between recall, \ac{FDR}, timing accuracy, and 
location reliability. In each subsection we systematically vary key 
hyperparameters (the choice of training datasets and detection probability 
thresholds for P- and S-phase detection) and analyze their influence on the 
performance metrics such as recall, timing accuracy, and consistency with the 
reference OGS catalog. Together, these results characterize the trade-offs 
between sensitivity and specificity inherent in automated seismic monitoring 
systems. Throughout this chapter, all comparisons are made against the manually 
reviewed OGS reference catalog.

\section{Picker}
The picker module is the first critical step in our seismic event detection 
pipeline. Its role is to identify P- and S-wave arrivals in continuous waveform 
data, providing the foundation upon which association, location, and magnitude 
estimation depend. In this section, we present the results of the evaluation of 
the picker module as a function of different training datasets and detection
thresholds. We focus on the comparative behavior of PhaseNet and EQTransformer,
two state-of-the-art deep learning models for P- and S-phase detection. These 
models have been trained on different datasets (INSTANCE, STEAD, SCEDC and the
original training set) to assess how training data composition influences
generalization to the OGS seismic network. We report the pick recall, as well
as the pick time residuals between our results and the OGS catalog as a
function of these parameters.

Two main questions guide this analysis:
\begin{enumerate}
  \item How sensitive is picker performance to the choice of training dataset, 
  particularly with respect to domain mismatch between the training data and 
  the OGS evaluation region?
  \item How does the detection threshold mediate the trade-off between 
  recall (recovering as many catalog picks as possible) and \ac{FDR} 
  (limiting the number of spurious or poorly timed picks)?
\end{enumerate}

\subsection{Influence of Training Datasets}
Table~\ref{tab:picker_recall_datasets} offers a comparative overview of
picker performance when trained on different datasets. The OGS reference
catalog contains 7,129 manually reviewed picks (4,126~P and 3,003~S) over the
93-day evaluation period, recorded across 266~stations in 18~networks. Note 
that small discrepancies in the total pick counts across confusion matrices 
arise from the bipartite matching procedure, which evaluates each model's 
output independently against the catalog. The pick-level confusion matrices 
(Table~\ref{tab:picker_recall_datasets}) quantify each model's ability to 
recover these catalog picks, while Table~\ref{tab:picker_recall_scores} 
summarizes the per-phase and combined recall scores.

\subsubsection{Confusion Matrix Analysis}
The confusion matrices in Table~\ref{tab:picker_recall_datasets} display, for
each model--dataset pair, the number of matched (MH), missed (MS), and proposed
(PS) picks, decomposed by P- and S-phase. Each cell encodes the complete
detection profile of the corresponding configuration.

\paragraph{INSTANCE-trained models.}
PhaseNet (INSTANCE) matches 3,932 of 4,126~P picks (recall~0.953, only 
190~$MS_P$ and 4~$\epsilon_{PS}$) and 2,757 of 3,003~S picks (recall~0.918, 
236~$MS_S$ and 10~$\epsilon_{SP}$), achieving a combined recall of 0.938. The 
confusion matrix shows that its proposed picks picks volume is 1,972,866, 
indicating a large but manageable background of candidate detections from which 
downstream association modules must select. EQTransformer (INSTANCE) attains 
marginally higher figures on P-phase (3,965~$MH_P$, recall~0.961, 155~$MS_P$ 
and 6~$\epsilon_{PS}$) and identical S-phase recall (2,764~$MH_S$, 
recall~0.920, 236~$MS_S$ and 3~$\epsilon_{SP}$), yielding a combined recall of 
0.944. Its proposed pick count is lower at 1,364,149, roughly $31\%$ fewer than 
PhaseNet (INSTANCE), which may reflect EQTransformer's more conservative 
internal detection logic. These two INSTANCE-trained models detect 
substantially more catalog picks than any other training-dataset variant.

\paragraph{SCEDC-trained models.}
PhaseNet (SCEDC) matches 3,679~P picks (recall~0.892, 385~$MS_P$ and 
62~$\epsilon_{PS}$) and 2,471~S picks (recall~0.823, 433~$MS_S$ and 
99~$\epsilon_{SP}$), for a combined recall of 0.863. EQTransformer (SCEDC) 
performs better with 3,921~$MH_P$ (recall~0.95) and 2,682~$MH_S$ 
(recall~0.893), achieving a combined recall of 0.926. However, the confusion 
matrices reveal the critical weakness of the SCEDC training set: PhaseNet 
(SCEDC) generates 23.1 million proposed picks and EQTransformer (SCEDC) nearly 
30 million, roughly 12$\times$ to 22$\times$ more than the INSTANCE-trained 
models. This massive over-detection inflates the apparent recall while burying
matched picks in an overwhelming volume of proposed picks, rendering the output 
impractical for downstream processing without aggressive pre-filtering. The 
SCEDC dataset, composed of Southern California seismicity, introduces a severe 
domain mismatch with the Alpine-Dinaric tectonic setting of the OGS network, 
manifesting as an inability to discriminate noise from signal in this context.

\paragraph{Original training sets.}
PhaseNet (Original) shows a distinctive asymmetric profile visible in its 
confusion matrix: P recall is high at 0.956 (3,943~$MH_P$, only 175~$MS_P$ and
8~$\epsilon_{PS}$), but S recall drops to 0.858 (2,577~$MH_S$, 402~$MS_S$ and 
24~$\epsilon_{SP}$). Interestingly, its proposed picks volume is 10.4~million, 
substantially larger than the INSTANCE variant, suggesting that the original 
model generates more spurious detections in addition to missing true 
S-arrivals. EQTransformer (Original) exhibits uniformly reduced performance on 
both phases (P: 0.864, 3,563~$MH_P$, 558~$MS_P$ and 5~$\epsilon_{PS}$; S: 
0.843, 2,532~$MH_S$, 467~$MS_S$ and 4~$\epsilon_{SP}$; combined: 0.855) but 
with a relatively modest proposed picks volume of 651,963, indicating a 
conservative detection strategy that sacrifices sensitivity.

\paragraph{STEAD-trained models.}
The STEAD-trained pickers perform worst for this network. PhaseNet (STEAD) 
recovers only 2,929~P picks (recall~0.71, 1,181~$MS_P$ and 16~$\epsilon_{PS}$) 
and 2,424~S picks (recall~0.807, 564~$MS_S$ and 15~$\epsilon_{SP}$), for a
combined recall of 0.75. The confusion matrix reveals the lowest proposed picks 
volume of any PhaseNet variant (978,247), confirming that STEAD training 
produces an under-sensitive model that misses a substantial fraction of real 
arrivals while also generating fewer candidates overall. EQTransformer (STEAD) 
fares somewhat better (P: 0.841, 3,470~$MH_P$, 651~$MS_P$, and 
5~$\epsilon_{PS}$; S: 0.825, 2,478~$MH_S$, 520~$MS_S$, and 5~$\epsilon_{SP}$; 
combined: 0.834) with 677,327~proposed picks, but remains well below the 
INSTANCE-trained models.\\

The INSTANCE dataset includes stations and events from Northeastern Italy
(including OGS stations), which likely contributes to the strong domain
alignment. The STEAD dataset, composed predominantly of global teleseismic and 
regional recordings with different noise characteristics, produces the least 
sensitive pickers for this network. Table~\ref{tab:picker_recall_scores}
summarizes the recall scores, confirming the INSTANCE advantage across both 
architectures.

\subsubsection{Pick Time Residual Analysis}
The pick time residual histograms in Table~\ref{tab:picker_timeres_datasets}
quantify the timing accuracy of each model by showing the distribution of
$\delta t = t_{\text{detected}} - t_{\text{catalog}}$ for all matched picks,
separately for P- and S-phases. These histograms reveal both the central
tendency (systematic bias) and the spread (precision) of each model's timing.

\paragraph{PhaseNet (INSTANCE).}
The P-phase residual histogram is sharply peaked near zero with a slight
negative bias: mean~$= -0.021$\,s, MAE~$= 0.066$\,s, RMSE~$= 0.107$\,s. The 
S-phase histogram is broader with a positive bias: mean~$= 0.054$\,s,
MAE~$= 0.116$\,s, RMSE~$= 0.16$\,s. The larger S-phase scatter reflects the 
intrinsic difficulty of timing S-arrivals, which emerge from the P-coda and are
typically less impulsive. The standard deviation of 0.151\,s for S-picks
compared to 0.105\,s for P-picks quantifies this phase-dependent uncertainty 
difference.

\paragraph{EQTransformer (INSTANCE).}
The P-phase residual histogram shows near-zero bias (mean~$= 0.003$\,s) but
slightly wider spread than PhaseNet: MAE~$= 0.075$\,s, RMSE~$= 0.114$\,s.
The S-phase residuals exhibit a more pronounced positive bias 
(mean~$=0.091$\,s, MAE~$= 0.134$\,s, RMSE~$= 0.179$\,s). Despite 
EQTransformer's marginally higher overall recall (0.944 vs.\ 0.938), PhaseNet's 
superior timing precision, 12\% lower P-phase MAE and 16\% lower S-phase MAE, 
makes it the more suitable choice for operational use, where accurate pick 
timing is crucial for successful association and location.

\paragraph{Non-INSTANCE models.}
Among the non-INSTANCE models, the time residual histograms reveal
characteristic biases linked to domain mismatch:
\begin{itemize}
  \item \textbf{PhaseNet (Original):} exhibits the largest systematic S-phase
    bias of any variant (mean~$=0.146$\,s, MAE~$=0.156$\,s, RMSE~$=0.200$\,s), 
    while its P-phase shows a positive bias of $0.046$\,s (MAE~$= 0.081$\,s).
  \item \textbf{PhaseNet (STEAD):} produces P residuals comparable to INSTANCE 
    (mean~$=-0.018$\,s, MAE~$=0.065$\,s) but at much lower recall (0.710 vs. 
    0.953), so the narrow histogram is partly an artifact of detecting only the 
    clearest arrivals. S-phase bias is moderate ($0.098$\,s, MAE~$=0.128$\,s).
  \item \textbf{EQTransformer (SCEDC):} shows a large S-phase bias
    (mean~$=0.131$\,s, MAE~$=0.156$\,s, RMSE~$=0.201$\,s).
  \item \textbf{EQTransformer (Original):} achieves the smallest P-phase bias 
    among EQTransformer variants (mean~$=-0.013$\,s, MAE~$= 0.071$\,s).
\end{itemize}

The timing advantage of INSTANCE-trained models is operationally significant: 
downstream association and location modules rely on sub-sample precision, so 
even differences of $\sim$0.02\,s in MAE can affect event characterization 
quality.

Based on these results, PhaseNet trained on INSTANCE is selected as the picker 
module for all subsequent analyses: it combines the second-highest overall 
recall (0.938) with the best timing accuracy across both P and S phases, making 
it the most suitable configuration for the OGS network.

\begin{table}
  \centering
  \caption{Recall scores for PhaseNet and EQTransformer models trained on 
  different data sets (detection threshold: prob $\geq$ 0.1).}
  \label{tab:picker_recall_scores}
  \begin{tabular}{c|c|c|c|c|c|c|c|c|c|}
    \setlength\heavyrulewidth{2pt} % does it have to be soooo wide?
    \setlength\extrarowheight{2pt}
    \setlength\tabcolsep{0pt}
    \multirow{2}{*}{Recall} & \multicolumn{3}{c|}{PhaseNet} & \multicolumn{3}{c|}{EQTransformer} \\
    & P & S & Tot & P & S & Tot \\
    \hline
    \textbf{INSTANCE} & \textbf{0.953} & \textbf{0.918} & \textbf{0.938} & \textbf{0.961} & \textbf{0.920} & \textbf{0.944} \\
    SCEDC             & 0.892 & 0.823 & 0.863 & 0.950 & 0.893 & 0.926 \\
    Original          & 0.956 & 0.858 & 0.915 & 0.864 & 0.843 & 0.855 \\
    STEAD             & 0.710 & 0.807 & 0.750 & 0.841 & 0.825 & 0.834 \\
    \hline
  \end{tabular}
\end{table}

\newpage
\begin{table}[H]
  \centering
  \begin{tabular}{c|c|c|}
    & PhaseNet & EQTransformer \\\hline
    Original & \includegraphics[width=0.35\textwidth]{OGSCatalog/OGSPhaseNet_ORIGINAL/img/.all_SeisBenchPicker_PicksConfMtx.pdf} & \includegraphics[width=0.35\textwidth]{OGSCatalog/OGSEQTransformer_ORIGINAL/img/.all_SeisBenchPicker_PicksConfMtx.pdf} \\\hline
    INSTANCE & \includegraphics[width=0.35\textwidth]{OGSCatalog/GaMMA0.1/img/.all_SeisBenchPicker_PicksConfMtx.pdf} & \includegraphics[width=0.35\textwidth]{OGSCatalog/OGSEQTransformer_INSTANCE/img/.all_SeisBenchPicker_PicksConfMtx.pdf} \\\hline
    STEAD & \includegraphics[width=0.35\textwidth]{OGSCatalog/OGSPhaseNet_STEAD/img/.all_SeisBenchPicker_PicksConfMtx.pdf} & \includegraphics[width=0.35\textwidth]{OGSCatalog/OGSEQTransformer_STEAD/img/.all_SeisBenchPicker_PicksConfMtx.pdf}\\\hline
    SCEDC & \includegraphics[width=0.35\textwidth]{OGSCatalog/OGSPhaseNet_SCEDC/img/.all_SeisBenchPicker_PicksConfMtx.pdf} & \includegraphics[width=0.35\textwidth]{OGSCatalog/OGSEQTransformer_SCEDC/img/.all_SeisBenchPicker_PicksConfMtx.pdf} \\\hline
  \end{tabular}
  \caption{Pick-level recall for PhaseNet and EQTransformer trained on four
  distinct datasets (detection threshold: prob $\geq$ 0.1).}
  \label{tab:picker_recall_datasets}
\end{table}

\newpage
\subsubsection{Pick Time Difference}
The pick time residual histograms in Table~\ref{tab:picker_timeres_datasets}
present the distribution of timing differences
($\delta t = t_{\text{detected}} - t_{\text{catalog}}$) for each model and
training dataset combination. Positive residuals indicate that the automated
pick arrives later than the manual catalog time; negative residuals indicate
earlier picks. The width of each histogram reflects the precision of the model, 
while its center quantifies systematic bias.
\begin{table}[H]
  \centering
  \begin{tabular}{c|c|c|}
    & PhaseNet & EQTransformer \\\hline
    Original & \includegraphics[width=0.4\textwidth]{OGSCatalog/OGSPhaseNet_ORIGINAL/img/.all_SeisBenchPicker_PicksTimeDiff.pdf} & \includegraphics[width=0.4\textwidth]{OGSCatalog/OGSEQTransformer_ORIGINAL/img/.all_SeisBenchPicker_PicksTimeDiff.pdf} \\\hline
    INSTANCE & \includegraphics[width=0.4\textwidth]{OGSCatalog/GaMMA0.1/img/.all_SeisBenchPicker_PicksTimeDiff.pdf} & \includegraphics[width=0.4\textwidth]{OGSCatalog/OGSEQTransformer_INSTANCE/img/.all_SeisBenchPicker_PicksTimeDiff.pdf} \\\hline
    STEAD & \includegraphics[width=0.4\textwidth]{OGSCatalog/OGSPhaseNet_STEAD/img/.all_SeisBenchPicker_PicksTimeDiff.pdf} & \includegraphics[width=0.4\textwidth]{OGSCatalog/OGSEQTransformer_STEAD/img/.all_SeisBenchPicker_PicksTimeDiff.pdf} \\\hline
    SCEDC & \includegraphics[width=0.4\textwidth]{OGSCatalog/OGSPhaseNet_SCEDC/img/.all_SeisBenchPicker_PicksTimeDiff.pdf} & \includegraphics[width=0.4\textwidth]{OGSCatalog/OGSEQTransformer_SCEDC/img/.all_SeisBenchPicker_PicksTimeDiff.pdf} \\\hline
  \end{tabular}
  \caption{Pick time residuals ($\delta t$, detected minus catalog times) for 
  PhaseNet (INSTANCE) and EQTransformer trained on different datasets
  (detection threshold: prob $\geq$ 0.1).}
  \label{tab:picker_timeres_datasets}
\end{table}
\newpage
\subsection{Influence of Detection Thresholds}
Deep learning phase pickers apply a probability threshold to accept or reject a
candidate phase detection. Here we systematically evaluate its effect on
PhaseNet (INSTANCE) and EQTransformer (INSTANCE) across three settings: 0.1,
0.2, and 0.3. The left columns of Tables~\ref{tab:phasenet_thresholds}
and~\ref{tab:eqt_thresholds} show the confusion matrices at each threshold,
while the right columns display the corresponding pick time residual
histograms.

\subsubsection{PhaseNet (INSTANCE) Threshold Sweep}
For PhaseNet (INSTANCE) (Table~\ref{tab:phasenet_thresholds}), raising the
threshold from 0.1 to 0.3 reduces total recall from 0.938 to 0.897. Examining
the confusion matrices row-by-row:
\begin{itemize}
  \item \textbf{Threshold 0.1:} 3,932~$MH_P$ (recall~0.953), 2,757~$MH_S$ 
    (recall~0.918), combined recall~0.938, with 1,972,866~PS.
  \item \textbf{Threshold 0.2:} 3,887~$MH_P$ (recall~0.942), 2,669~$MH_S$ 
    (recall~0.889), combined recall~0.920, with 790,480~PS, a 60\% reduction 
    from threshold 0.1.
  \item \textbf{Threshold 0.3:} 3,835~$MH_P$ (recall~0.929), 2,561~$MH_S$ 
    (recall~0.853), combined recall~0.897, with 398,815~PS, a factor of 
    $\sim$5$\times$ reduction from threshold 0.1.
\end{itemize}
The S-phase recall declines more steeply than P-phase recall (6.6 percentage 
points vs.\ 2.4 percentage points from 0.1 to 0.3), confirming that S-wave 
detections are more threshold-sensitive due to their generally lower confidence 
scores.

The time residual histograms in the right column of
Table~\ref{tab:phasenet_thresholds} show remarkable stability across
thresholds. The P-phase MAE remains virtually constant at $\sim$0.065\,s 
(0.066\,s at 0.1, 0.065\,s at 0.2, 0.065\,s at 0.3), with the mean bias stable 
at $-0.021$\,s and median at $-0.020$\,s. The S-phase MAE improves marginally 
from 0.116\,s to 0.110\,s as the threshold rises, with RMSE decreasing from 
0.160\,s to 0.153\,s. This confirms that higher thresholds primarily filter out 
poorly timed or spurious picks without degrading timing quality for the 
retained detections. The histogram shapes narrow slightly at higher thresholds,
indicating that low-confidence picks, which tend to have larger timing errors, 
are preferentially removed.

This threshold-dependent trade-off is critical for operational monitoring: 
lower thresholds enhance sensitivity to weak events, whereas higher thresholds 
favor accurate timing for downstream association algorithms.
\begin{table}[H]
  \centering
  \begin{tabular}{c|c|c|c|}
    & Confusion Matrix & Time Residual Histogram \\\hline
    0.1 & \includegraphics[width=0.35\textwidth]{OGSCatalog/GaMMA0.1/img/.all_SeisBenchPicker_PicksConfMtx.pdf} & \includegraphics[width=0.55\textwidth]{OGSCatalog/GaMMA0.1/img/.all_SeisBenchPicker_PicksTimeDiff.pdf}\\\hline
    0.2 & \includegraphics[width=0.35\textwidth]{OGSCatalog/GaMMA0.2/img/.all_SeisBenchPicker_PicksConfMtx.pdf} & \includegraphics[width=0.55\textwidth]{OGSCatalog/GaMMA0.2/img/.all_SeisBenchPicker_PicksTimeDiff.pdf}\\\hline
    0.3 & \includegraphics[width=0.35\textwidth]{OGSCatalog/GaMMA0.3/img/.all_SeisBenchPicker_PicksConfMtx.pdf} & \includegraphics[width=0.55\textwidth]{OGSCatalog/GaMMA0.3/img/.all_SeisBenchPicker_PicksTimeDiff.pdf}\\\hline
  \end{tabular}
  \caption{[PhaseNet (INSTANCE)] Picker recall and pick time difference as a 
  function of P and S detection thresholds (0.1, 0.2, 0.3).}
  \label{tab:phasenet_thresholds}
\end{table}

\newpage
\subsubsection{EQTransformer (INSTANCE) Threshold Sweep}
Table~\ref{tab:eqt_thresholds} presents the corresponding analysis for
EQTransformer (INSTANCE). The same general pattern holds, with recall
decreasing monotonically as the threshold rises:
\begin{itemize}
  \item \textbf{Threshold 0.1:} 3,965~$MH_P$ (recall~0.961), 2,764~$MH_S$ 
  (recall~0.920), combined 0.944, with 1,364,149~PS.
  \item \textbf{Threshold 0.2:} 3,934~$MH_P$ (recall~0.953), 2,687~$MH_S$ 
  (recall~0.895), combined 0.929, with 632,974~PS.
  \item \textbf{Threshold 0.3:} 3,888~$MH_P$ (recall~0.942), 2,619~$MH_S$ 
  (recall~0.872), combined 0.913, with 366,570~PS.
\end{itemize}
The recall decline from 0.1 to 0.3 is slightly smaller than PhaseNet's, with 
the S-phase experiencing a steeper drop than the P-phase, mirroring the 
phase-dependent sensitivity observed with PhaseNet.

EQTransformer's P-phase timing is nearly threshold-invariant: MAE={0.075, 
0.074, 0.074}\,s across thresholds \{0.1, 0.2, 0.3\}, with mean bias stable at 
$+$0.003\,s. The S-phase MAE improves slightly from 0.134\,s to 0.129\,s, with 
RMSE decreasing from 0.179\,s to 0.174\,s. Compared with PhaseNet (INSTANCE) at 
each threshold level, EQTransformer retains marginally higher recall (by 
0.5--1.4 percentage points) but consistently exhibits larger timing residuals, 
approximately 15\% higher P-phase MAE and 18\% higher S-phase MAE. This 
confirms the fundamental architectural trade-off: EQTransformer's attention 
mechanism achieves slightly broader detection sensitivity at the cost of 
reduced timing precision.
\begin{table}[H]
  \centering
  \begin{tabular}{c|c|c|}
    & Confusion Matrix & Time Residual Histogram \\\hline
    0.1 & \includegraphics[width=0.35\textwidth]{OGSCatalog/OGSEQTransformer_INSTANCE/img/.all_SeisBenchPicker_PicksConfMtx.pdf} & \includegraphics[width=0.55\textwidth]{OGSCatalog/OGSEQTransformer_INSTANCE/img/.all_SeisBenchPicker_PicksTimeDiff.pdf}\\\hline
    0.2 & \includegraphics[width=0.35\textwidth]{OGSCatalog/OGSEQTransformer_INSTANCE_TP0.2S0.2/img/.all_SeisBenchPicker_PicksConfMtx.pdf} & \includegraphics[width=0.55\textwidth]{OGSCatalog/OGSEQTransformer_INSTANCE_TP0.2S0.2/img/.all_SeisBenchPicker_PicksTimeDiff.pdf}\\\hline
    0.3 & \includegraphics[width=0.35\textwidth]{OGSCatalog/OGSEQTransformer_INSTANCE_TP0.3S0.3/img/.all_SeisBenchPicker_PicksConfMtx.pdf} & \includegraphics[width=0.55\textwidth]{OGSCatalog/OGSEQTransformer_INSTANCE_TP0.3S0.3/img/.all_SeisBenchPicker_PicksTimeDiff.pdf}\\\hline
  \end{tabular}
  \caption{[EQTransformer (INSTANCE)] Picker recall and pick time difference as 
  a function of P and S detection thresholds (0.1, 0.2, 0.3).}
  \label{tab:eqt_thresholds}
\end{table}

\newpage
\section{Associator}
Having established that PhaseNet trained on INSTANCE yields the best balance of 
recall and timing precision, the next question is how effectively association 
algorithms group its output picks into coherent seismic events. The associator 
module is the second critical step in our seismic event detection pipeline: it 
groups detected phase picks into coherent seismic events by enforcing 
spatiotemporal consistency among picks to form events and intentionally 
discards isolated or inconsistent picks that cannot be reliably attributed to 
any event. Its effectiveness determines how many of the 380 OGS catalog events 
the pipeline can recover. After fixing the picker to PhaseNet (INSTANCE), we 
evaluate two associators, GaMMA and PyOcto, as a function of pick detection 
thresholds (0.1, 0.2, 0.3). We report:
\begin{enumerate}
  \item Pick-level recall after association (how many matched picks survive),
  \item Event-level recall (how many of the 380 catalog events are recovered),
  \item Event time residuals, and
  \item Spatial residuals (epicentral distance and depth differences).
\end{enumerate}
A key distinction emerges: while the raw picker detects $\sim$6,689
true-positive picks (3,932 P $+$ 2,757 S at threshold 0.1), the associators 
retain only the subset that forms consistent event hypotheses, necessarily 
reducing pick-level recall but filtering the massive proposed picks volume 
($\sim$2 million at threshold 0.1) down to manageable event counts of 
$\sim$1,000--1,700 events.

\subsection{GaMMA}
The GaMMA associator is designed to probabilistically associate picks into
events by jointly modeling pick timing uncertainties and tentative source
locations. Tables~\ref{tab:gamma_picks} and~\ref{tab:gamma_events} present the 
pick-level and event-level results after GaMMA processing, respectively.

\subsubsection{Pick-Level Results}
The confusion matrices in Table~\ref{tab:gamma_picks} reveal the substantial
filtering that GaMMA imposes on the raw PhaseNet output:
\begin{itemize}
  \item \textbf{Threshold 0.1:} GaMMA retains 2,949 of 4,126~P picks (P recall 
    0.715) and 2,081 of 3,003~S picks (S recall 0.693), for a combined pick 
    recall of 0.706, substantially lower than the raw PhaseNet (INSTANCE) 
    output (0.938). This indicates that GaMMA discards $\sim$25\% of 
    true-positive picks during association. The 32,946 proposed picks that 
    survive association represent a dramatic reduction from the $\sim$2~million 
    raw PhaseNet proposals, filtering the pick set by a factor of 
    $\sim$60$\times$.
  \item \textbf{Threshold 0.2:} Pick recall decreases to 0.696 (P: 0.708,
    2,921~$MH_P$; S: 0.680, 2,041~$MH_S$), with proposed picks reduced to
    25,168.
  \item \textbf{Threshold 0.3:} Pick recall falls further to 0.665 (P: 0.683, 
    2,820~$MH_P$; S: 0.640, 1,921~$MH_S$), losing an additional 289 matched 
    picks relative to threshold 0.1. Proposed picks decrease to 19,538.
\end{itemize}

The time residual histograms in the right column of Table~\ref{tab:gamma_picks}
show that GaMMA preferentially retains well-timed arrivals. The P-phase MAE
improves from 0.064\,s (at 0.1) to 0.062\,s (at 0.3), with the mean bias stable 
at approximately $-0.017$\,s to $-0.020$\,s. The S-phase residuals remain 
consistent: MAE ranges from 0.111\,s (at 0.1) to 0.107\,s (at 0.3), with a 
positive bias of approximately $0.052$\,s. These values are comparable to or 
slightly better than the raw PhaseNet residuals, confirming that GaMMA's 
probabilistic filtering selectively discards poorly timed picks.

\subsubsection{Event-Level Results}
The event-level confusion matrices in Table~\ref{tab:gamma_events} show that
GaMMA achieves high recall despite the aggressive pick-level filtering:
\begin{itemize}
  \item \textbf{Threshold 0.1:} 365 of 380 catalog events are matched
    (recall~0.961), with only 15~$MS_E$. GaMMA proposes 1,060 additional events 
    absent from the OGS catalog (event \ac{FDR}~0.744).
  \item \textbf{Threshold 0.2:} 363~$MH_E$ (recall~0.955), 17~$MS_E$, 
    851~$PS_E$ (\ac{FDR}~0.701).
  \item \textbf{Threshold 0.3:} 360~$MH_E$ (recall~0.947), 20~$MS_E$, 
    688~$PS_E$ (\ac{FDR}~0.656).
\end{itemize}

The event time residual histograms in the right column of
Table~\ref{tab:gamma_events} reveal a consistent positive bias in GaMMA's 
origin time estimates: the mean origin time difference is $0.484$\,s at 
threshold 0.1 (MAE~$= 0.536$\,s, RMSE~$= 0.627$\,s, median~$= 0.493$\,s), 
improving slightly to $0.450$\,s at threshold 0.3 (MAE~$= 0.501$\,s, 
RMSE~$= 0.581$\,s). This systematic positive bias means that GaMMA's estimated 
origin times are consistently later than the OGS catalog values by 
approximately 0.5\,s, which is a consequence of GaMMA's internal origin-time 
estimation methodology that differs from the manual processing workflow.

\subsubsection{Spatial Residuals}
The epicentral distance and depth difference histograms (landscape table
following Table~\ref{tab:gamma_events}) characterize GaMMA's preliminary
location accuracy. At threshold 0.1, the mean epicentral distance difference is 
1.56\,km (median 1.10\,km, 90th percentile 3.30\,km), indicating that most 
GaMMA-located events fall within 1--3\,km of their OGS catalog positions. The 
depth residuals show a slight negative bias (mean~$= -0.78$\,km, 
median~$= -0.75$\,km), with a standard deviation of 3.38\,km and 90th 
percentile ranging from $-4.87$\,km to $+2.89$\,km. At higher thresholds, the 
spatial residuals improve modestly (median epicentral distance 1.01\,km at 
threshold 0.3), as cleaner pick sets yield more constrained events.

\begin{table}[H]
  \centering
  \begin{tabular}{c|c|c|}
    & Confusion Matrix & Time Residual Histogram \\
    \hline
    0.1 & \includegraphics[width=0.35\textwidth]{OGSCatalog/GaMMA0.1/img/.all_OGSPickStatQC_PicksConfMtx.pdf} & \includegraphics[width=0.55\textwidth]{OGSCatalog/GaMMA0.1/img/.all_OGSPickStatQC_PicksTimeDiff.pdf} \\\hline
    0.2 & \includegraphics[width=0.35\textwidth]{OGSCatalog/GaMMA0.2/img/.all_OGSPickStatQC_PicksConfMtx.pdf} & \includegraphics[width=0.55\textwidth]{OGSCatalog/GaMMA0.2/img/.all_OGSPickStatQC_PicksTimeDiff.pdf} \\\hline
    0.3 & \includegraphics[width=0.35\textwidth]{OGSCatalog/GaMMA0.3/img/.all_OGSPickStatQC_PicksConfMtx.pdf} & \includegraphics[width=0.55\textwidth]{OGSCatalog/GaMMA0.3/img/.all_OGSPickStatQC_PicksTimeDiff.pdf} \\\hline
  \end{tabular}
  \caption{[PhaseNet (INSTANCE) | GaMMA] Pick recall and pick time difference 
  as a function of P and S thresholds.}
  \label{tab:gamma_picks}
\end{table}

\newpage
\begin{table}[H]
  \centering
  \begin{tabular}{c|c|c|}
    & Confusion Matrix & Time Residual Histogram \\\hline
    0.1 & \includegraphics[width=0.35\textwidth]{OGSCatalog/GaMMA0.1/img/.all_OGSPickStatQC_EventsConfMtx.pdf} & \includegraphics[width=0.55\textwidth]{OGSCatalog/GaMMA0.1/img/.all_OGSPickStatQC_EventsTimeDiff.pdf} \\\hline
    0.2 & \includegraphics[width=0.35\textwidth]{OGSCatalog/GaMMA0.2/img/.all_OGSPickStatQC_EventsConfMtx.pdf} & \includegraphics[width=0.55\textwidth]{OGSCatalog/GaMMA0.2/img/.all_OGSPickStatQC_EventsTimeDiff.pdf} \\\hline
    0.3 & \includegraphics[width=0.35\textwidth]{OGSCatalog/GaMMA0.3/img/.all_OGSPickStatQC_EventsConfMtx.pdf} & \includegraphics[width=0.55\textwidth]{OGSCatalog/GaMMA0.3/img/.all_OGSPickStatQC_EventsTimeDiff.pdf} \\\hline
  \end{tabular}
  \caption{[PhaseNet (INSTANCE) | GaMMA] Event recall and event time difference 
  as a function of P and S thresholds.}
  \label{tab:gamma_events}
\end{table}

\newpage
\begin{landscape}
  \begin{table}[H]
    \centering
    \begin{tabular}{c|c|c|}
      & Epicentral Distance Difference & Depth Difference \\\hline
      0.1 & \includegraphics[width=0.55\textwidth]{OGSCatalog/GaMMA0.1/img/.all_OGSPickStatQC_EpiDistDiff.pdf} & \includegraphics[width=0.55\textwidth]{OGSCatalog/GaMMA0.1/img/.all_OGSPickStatQC_DepthDiff.pdf} \\\hline
      0.2 & \includegraphics[width=0.55\textwidth]{OGSCatalog/GaMMA0.2/img/.all_OGSPickStatQC_EpiDistDiff.pdf} & \includegraphics[width=0.55\textwidth]{OGSCatalog/GaMMA0.2/img/.all_OGSPickStatQC_DepthDiff.pdf} \\\hline
      0.3 & \includegraphics[width=0.55\textwidth]{OGSCatalog/GaMMA0.3/img/.all_OGSPickStatQC_EpiDistDiff.pdf} & \includegraphics[width=0.55\textwidth]{OGSCatalog/GaMMA0.3/img/.all_OGSPickStatQC_DepthDiff.pdf} \\\hline
    \end{tabular}
    \caption{[PhaseNet (INSTANCE) | GaMMA] Epicentral distance and depth
    differences as a function of P and S thresholds.}
  \end{table}
\end{landscape}

\newpage
\subsection{PyOcto}
PyOcto implements a velocity-model-aware association strategy, clustering picks
in both time and space using travel-time predictions to identify candidate
events. Tables~\ref{tab:pyocto_picks} and~\ref{tab:pyocto_events} present the
pick-level and event-level results after PyOcto processing, evaluated at the
same thresholds as GaMMA for direct comparison.

\subsubsection{Pick-Level Results}
The confusion matrices in Table~\ref{tab:pyocto_picks} demonstrate that
PyOcto's velocity-model-aware approach is substantially less aggressive in
filtering true-positive picks compared to GaMMA:
\begin{itemize}
  \item \textbf{Threshold 0.1:} PyOcto retains 3,875 of 4,126~P picks (P 
    recall~0.939) and 2,589 of 3,003~S picks (S recall~0.862), achieving a 
    combined pick recall of 0.907, substantially higher than GaMMA's 0.706. 
    PyOcto preserves 94\% of $MH_P$, compared with GaMMA's 72\%, demonstrating 
    that its spatial-consistency criterion is less aggressive than GaMMA's 
    probabilistic filtering. However, PyOcto also retains 51,485 proposed 
    picks, roughly 56\% more than GaMMA's 32,946, reflecting the trade-off 
    between pick retention and selectivity.
  \item \textbf{Threshold 0.2:} 3,813~$MH_P$ (recall~0.924), 2,530~$MH_S$
    (recall~0.842), combined recall~0.89, with 37,460 proposed picks.
  \item \textbf{Threshold 0.3:} 3,759~$MH_P$ (recall~0.911), 2,435~$MH_S$
    (recall~0.811), combined recall~0.869, with 28,352 proposed picks, a 
    $\sim$45\% reduction from threshold 0.1.
\end{itemize}

The time residual histograms in the right column of
Table~\ref{tab:pyocto_picks} are nearly identical to the raw PhaseNet values:
P-phase MAE~$\approx 0.066$\,s (mean~$= -0.021$\,s), S-phase 
MAE~$\approx 0.115$\,s (mean~$=0.056$\,s) at threshold 0.1, improving slightly 
to P MAE~$= 0.064$\,s and S MAE~$= 0.110$\,s at threshold 0.3.

\subsubsection{Event-Level Results}
At the event level (Table~\ref{tab:pyocto_events}), PyOcto achieves identical 
recall to GaMMA at threshold 0.1: 365 of 380~events (recall~0.961), with 15 
missed events. The decline at higher thresholds is slightly steeper than 
GaMMA's:
\begin{itemize}
  \item \textbf{Threshold 0.2:} 361~$MH_E$ (recall~0.950), 19~$MS_E$,
    933~$PS_E$ (\ac{FDR}~0.721).
  \item \textbf{Threshold 0.3:} 356~$MH_E$ (recall~0.937), 24~$MS_E$,
    730~$PS_E$ (\ac{FDR}~0.672).
\end{itemize}
PyOcto proposes more events than GaMMA at each threshold (1,298 vs. 1,060 at 
0.1), yielding a higher event FDR of 0.781, which may include a larger fraction 
of genuine uncataloged micro-seismicity given PyOcto's stricter spatial 
constraints.

\paragraph{Cross-configuration FDR comparison.}
Across both associators and all thresholds, event-level \ac{FDR} remains high,
ranging from 0.656 (GaMMA, threshold~0.3) to 0.781 (PyOcto, threshold~0.1). At
the default threshold of 0.1, GaMMA produces roughly 2.9~$PS_E$ per matched
event, while PyOcto produces 3.6. Raising the threshold to 0.3 reduces these
ratios to $\sim$1.9 for GaMMA and $\sim$2.1 for PyOcto, demonstrating that
stricter detection thresholds provide a direct mechanism for \ac{FDR} reduction
at the cost of modest recall loss. However, these \ac{FDR} values should be
interpreted with caution: as noted in Section~\ref{chap:Methodology}, the
metric is computed against the OGS reference catalog, which is itself
incomplete below $M_L \approx 1$. A substantial fraction of the proposed events
may therefore represent genuine uncataloged micro-seismicity rather than false
detections, meaning the true operational false-positive rate is likely lower
than the nominal \ac{FDR} suggests.

The event time residual histograms reveal a contrasting systematic bias
compared to GaMMA: PyOcto's origin time estimates are consistently earlier than 
the OGS catalog, with mean~$= -0.421$\,s.

\subsubsection{Spatial Residuals}
The epicentral distance and depth difference histograms
(Table~\ref{tab:pyocto_spatial}) reveal PyOcto's spatial accuracy. At threshold 
0.1, the median epicentral distance difference is 1.29\,km (mean~1.65\,km, 90th 
percentile 3.32\,km), comparable to GaMMA's 1.10\,km. However, the depth 
residuals show markedly different behavior: PyOcto's mean depth difference is 
near zero ($-0.02$\,km), but with a much larger standard deviation of 5.03\,km 
compared to GaMMA's 3.38\,km. The 10th--90th percentile range spans from
$-7.89$\,km to $+5.02$\,km, substantially wider than GaMMA's range of 
$-4.87$\,km to $+2.89$\,km. This suggests that while PyOcto achieves unbiased 
depth estimates on average, its depth uncertainties are larger for individual 
events, a trade-off for retaining more events and picks.

\begin{table}[H]
  \centering
  \begin{tabular}{c|c|c|}
    & Confusion Matrix & Time Residual Histogram \\\hline
    0.1 & \includegraphics[width=0.35\textwidth]{OGSCatalog/PyOcto0.1/img/.all_OGSPickStatQC_PicksConfMtx.pdf} & \includegraphics[width=0.55\textwidth]{OGSCatalog/PyOcto0.1/img/.all_OGSPickStatQC_PicksTimeDiff.pdf} \\\hline
    0.2 & \includegraphics[width=0.35\textwidth]{OGSCatalog/PyOcto0.2/img/.all_OGSPickStatQC_PicksConfMtx.pdf} & \includegraphics[width=0.55\textwidth]{OGSCatalog/PyOcto0.2/img/.all_OGSPickStatQC_PicksTimeDiff.pdf} \\\hline
    0.3 & \includegraphics[width=0.35\textwidth]{OGSCatalog/PyOcto0.3/img/.all_OGSPickStatQC_PicksConfMtx.pdf} & \includegraphics[width=0.55\textwidth]{OGSCatalog/PyOcto0.3/img/.all_OGSPickStatQC_PicksTimeDiff.pdf} \\\hline
  \end{tabular}
  \caption{[PhaseNet (INSTANCE) | PyOcto] Pick recall and pick time difference 
  as a function of P and S thresholds.}
  \label{tab:pyocto_picks}
\end{table}

\begin{table}[H]
  \centering
  \begin{tabular}{c|c|c|}
    & Confusion Matrix & Time Residual Histogram \\\hline
    0.1 & \includegraphics[width=0.35\textwidth]{OGSCatalog/PyOcto0.1/img/.all_OGSPickStatQC_EventsConfMtx.pdf} & \includegraphics[width=0.55\textwidth]{OGSCatalog/PyOcto0.1/img/.all_OGSPickStatQC_EventsTimeDiff.pdf} \\\hline
    0.2 & \includegraphics[width=0.35\textwidth]{OGSCatalog/PyOcto0.2/img/.all_OGSPickStatQC_EventsConfMtx.pdf} & \includegraphics[width=0.55\textwidth]{OGSCatalog/PyOcto0.2/img/.all_OGSPickStatQC_EventsTimeDiff.pdf} \\\hline
    0.3 & \includegraphics[width=0.35\textwidth]{OGSCatalog/PyOcto0.3/img/.all_OGSPickStatQC_EventsConfMtx.pdf} & \includegraphics[width=0.55\textwidth]{OGSCatalog/PyOcto0.3/img/.all_OGSPickStatQC_EventsTimeDiff.pdf} \\\hline
  \end{tabular}
  \caption{[PhaseNet (INSTANCE) | PyOcto] Event recall and event time 
  difference as a function of P and S thresholds.}
  \label{tab:pyocto_events}
\end{table}

\newpage
\begin{landscape}
  \begin{table}[H]
    \centering
    \begin{tabular}{c|c|c|}
      & Epicentral Distance Difference & Depth Difference \\\hline
      0.1 & \includegraphics[width=0.55\textwidth]{OGSCatalog/PyOcto0.1/img/.all_OGSPickStatQC_EpiDistDiff.pdf} & \includegraphics[width=0.55\textwidth]{OGSCatalog/PyOcto0.1/img/.all_OGSPickStatQC_DepthDiff.pdf} \\\hline
      0.2 & \includegraphics[width=0.55\textwidth]{OGSCatalog/PyOcto0.2/img/.all_OGSPickStatQC_EpiDistDiff.pdf} & \includegraphics[width=0.55\textwidth]{OGSCatalog/PyOcto0.2/img/.all_OGSPickStatQC_DepthDiff.pdf} \\\hline
      0.3 & \includegraphics[width=0.55\textwidth]{OGSCatalog/PyOcto0.3/img/.all_OGSPickStatQC_EpiDistDiff.pdf} & \includegraphics[width=0.55\textwidth]{OGSCatalog/PyOcto0.3/img/.all_OGSPickStatQC_DepthDiff.pdf} \\\hline
    \end{tabular}
    \caption{[PhaseNet (INSTANCE) | PyOcto] Epicentral distance and depth
    differences as a function of P and S thresholds.}
    \label{tab:pyocto_spatial}
  \end{table}
\end{landscape}

\newpage
\section{NonLinLoc}
\label{sec:nonlinloc}
After evaluating the associator modules independently, we now examine the
combined performance of each associator followed by NonLinLoc (NonLinear
Location) as the location engine. NonLinLoc inverts associated picks for
earthquake hypocenters and origin times using a probabilistic non-linear
search. NonLinLoc is a widely used probabilistic location algorithm that can
utilize both 1D and 3D velocity models. We use a 1D layered velocity model for
all configurations to enable a controlled comparison of how each associator's
output survives the location stage.

\subsection{1D Velocity Model}
This stage reveals a striking divergence between the two associators. At the
association level, GaMMA and PyOcto achieve identical event recall (96.1\%,
365 of 380 events). After NonLinLoc location, however, GaMMA+NonLinLoc retains
only 280~$MH_E$ (recall~0.737), while PyOcto+NonLinLoc retains 366
events (recall~0.963). GaMMA loses 85 catalog events, nearly a quarter, during
location, whereas PyOcto actually recovers one additional event (366 vs. 365 at 
association), presumably because its pick sets yield more robust hypocentral 
solutions that occasionally match catalog events not previously associated.

\subsubsection{GaMMA}
The GaMMA+NonLinLoc 1D pathway shows a large recall drop relative to the
association stage. The confusion matrices in Table~\ref{tab:gamma_nll1d_events}
quantify this loss at each threshold:
\begin{itemize}
  \item \textbf{Threshold 0.1:} Only 280 of 380 catalog events are successfully 
    located (recall~0.737), down from 365 at association, a loss of 85 events (23\% relative drop). The 810~$PS_E$ that
    survive location (down from 1,060 at association) give an event \ac{FDR} of 
    0.743.
  \item \textbf{Threshold 0.2:} Recall decreases further to 0.716 (272~$MH_E$, 
    108~$MS_E$), with 669~$PS_E$ (\ac{FDR}~0.711).
  \item \textbf{Threshold 0.3:} Recall slightly recovers to 0.732 (278~$MH_E$, 
    102~$MS_E$), with 535~$PS_E$ (\ac{FDR}~0.658), as the cleaner pick set 
    partially compensates for reduced input volume.
\end{itemize}

The origin time residual histograms in the right column of
Table~\ref{tab:gamma_nll1d_events} show a marked improvement over the
association-stage estimates. The NonLinLoc inversion largely eliminates GaMMA's 
$0.5$\,s systematic bias: the mean origin time residual at threshold 0.1 is 
only $-0.062$\,s (MAE~$= 0.161$\,s, RMSE~$= 0.210$\,s, median~$= -0.066$\,s). 
At threshold 0.2, the residuals tighten further (MAE~$= 0.150$\,s, 
RMSE~$= 0.208$\,s), and at threshold 0.3 the MAE is 0.154\,s. The residual 
histograms are approximately Gaussian and centered near zero, confirming that 
NonLinLoc's probabilistic inversion effectively re-estimates origin times 
independent of the associator's preliminary values.

\paragraph{Post-location FDR.}
Notably, the NonLinLoc location stage does not substantially reduce event-level
\ac{FDR} for the GaMMA pathway. At threshold~0.1, \ac{FDR} changes only from
0.744 at association to 0.743 after location, indicating that NonLinLoc
removes matched and proposed events in comparable proportions. The same
behavior is observed at higher thresholds: \ac{FDR} is 0.711 at threshold~0.2
and 0.658 at threshold~0.3. Therefore, for GaMMA+NonLinLoc, the primary lever
for controlling post-location proposed event volume is the picker probability
threshold rather than the location step itself.

\paragraph{Spatial residuals.}
The epicentral distance and depth difference histograms
(Table~\ref{tab:gamma_nll1d_spatial}) demonstrate substantial improvement over 
the association-stage locations. At threshold 0.1:
\begin{itemize}
  \item \textbf{Epicentral distance:} median 0.67\,km (mean~1.11\,km,
    RMSE~1.75\,km), with 90\% of events located within 2.68\,km of the OGS 
    catalog position. This represents a $\sim$39\% improvement in median 
    epicentral accuracy compared to GaMMA's association-level estimates (median 
    1.10\,km).
  \item \textbf{Depth:} mean $0.50$\,km (median $0.72$\,km, std 2.75\,km), 
    with the interquartile range spanning $-0.40$\,km to $+1.99$\,km. The 
    positive bias indicates that NonLinLoc consistently places events slightly 
    deeper than the OGS catalog.
\end{itemize}
At threshold 0.3, the median epicentral distance remains stable (0.66\,km) 
while depth scatter is slightly reduced (std~2.71\,km), confirming that higher 
thresholds produce marginally cleaner location estimates.

\begin{table}[H]
  \centering
  \begin{tabular}{c|c|c|}
    & Confusion Matrix & Time Residual Histogram \\\hline
    0.1 & \includegraphics[width=0.35\textwidth]{OGSCatalog/GaMMA0.1/img/.all_OGSLocalMagnitude_EventsConfMtx.pdf} & \includegraphics[width=0.55\textwidth]{OGSCatalog/GaMMA0.1/img/.all_OGSLocalMagnitude_EventsTimeDiff.pdf} \\\hline
    0.2 & \includegraphics[width=0.35\textwidth]{OGSCatalog/GaMMA0.2/img/.all_OGSLocalMagnitude_EventsConfMtx.pdf} & \includegraphics[width=0.55\textwidth]{OGSCatalog/GaMMA0.2/img/.all_OGSLocalMagnitude_EventsTimeDiff.pdf} \\\hline
    0.3 & \includegraphics[width=0.35\textwidth]{OGSCatalog/GaMMA0.3/img/.all_OGSLocalMagnitude_EventsConfMtx.pdf} & \includegraphics[width=0.55\textwidth]{OGSCatalog/GaMMA0.3/img/.all_OGSLocalMagnitude_EventsTimeDiff.pdf} \\\hline
  \end{tabular}
  \caption{[PhaseNet (INSTANCE) | GaMMA | NonLinLoc 1D] Event detection and 
  origin time accuracy using 1D velocity model. Left column: event-level 
  confusion matrices comparing NonLinLoc-located events against the OGS 
  reference catalog. Right column: origin time residuals between the automated 
  workflow and OGS solutions.}
  \label{tab:gamma_nll1d_events}
\end{table}

\newpage
\begin{landscape}
  \begin{table}[H]
    \centering
    \begin{tabular}{c|c|c|}
      & Epicentral Distance Difference & Depth Difference \\\hline
      0.1 & \includegraphics[width=0.55\textwidth]{OGSCatalog/GaMMA0.1/img/.all_OGSLocalMagnitude_EpiDistDiff.pdf} & \includegraphics[width=0.55\textwidth]{OGSCatalog/GaMMA0.1/img/.all_OGSLocalMagnitude_DepthDiff.pdf} \\\hline
      0.2 & \includegraphics[width=0.55\textwidth]{OGSCatalog/GaMMA0.2/img/.all_OGSLocalMagnitude_EpiDistDiff.pdf} & \includegraphics[width=0.55\textwidth]{OGSCatalog/GaMMA0.2/img/.all_OGSLocalMagnitude_DepthDiff.pdf} \\\hline
      0.3 & \includegraphics[width=0.55\textwidth]{OGSCatalog/GaMMA0.3/img/.all_OGSLocalMagnitude_EpiDistDiff.pdf} & \includegraphics[width=0.55\textwidth]{OGSCatalog/GaMMA0.3/img/.all_OGSLocalMagnitude_DepthDiff.pdf} \\\hline
    \end{tabular}
    \caption{[PhaseNet (INSTANCE) | GaMMA | NonLinLoc 1D] Epicentral distance 
    and depth differences as a function of P and S thresholds.}
    \label{tab:gamma_nll1d_spatial}
  \end{table}
\end{landscape}

\subsubsection{PyOcto}
In contrast, the PyOcto+NonLinLoc 1D pathway
(Table~\ref{tab:pyocto_nll1d_events}) preserves nearly all associated events
through the location stage. The confusion matrices show:
\begin{itemize}
  \item \textbf{Threshold 0.1:} 366 of 380 catalog events are located
    (recall~0.963), with only 14~$MS_E$, essentially matching the 365 events 
    from association while recovering one additional event. The 1,304~$PS_E$ 
    (\ac{FDR}~0.781) suggest abundant new detections beyond the OGS catalog.
  \item \textbf{Threshold 0.2:} 359~$MH_E$ (recall~0.945), 21~$MS_E$, 
    947~$PS_E$ (\ac{FDR}~0.725).
  \item \textbf{Threshold 0.3:} 353~$MH_E$ (recall~0.929), 27~$MS_E$, 
    734~$PS_E$ (\ac{FDR}~0.675).
\end{itemize}
Even at the most restrictive threshold, PyOcto+NonLinLoc locates 73 more
catalog events than GaMMA+NonLinLoc at threshold 0.1 (353 vs. 280).\\

The origin time residual histograms show characteristics similar to the GaMMA 
pathway after NonLinLoc inversion: mean~$= -0.061$\,s at threshold 0.1 
(MAE~$=0.171$\,s, RMSE~$=0.236$\,s, median~$=-0.069$\,s). At threshold 0.3, the 
residuals tighten (MAE~$=0.145$\,s, RMSE~$=0.196$\,s), representing the best 
origin time accuracy across all configurations. The slight negative bias
($\sim -0.06$\,s) is common to both pathways after NonLinLoc processing,
suggesting a systematic difference between the automated and manual
origin-time estimation procedures.

\paragraph{Post-location FDR.}
Notably, the NonLinLoc location stage does not substantially reduce \ac{FDR}
for either pathway. GaMMA+NonLinLoc retains an \ac{FDR} of 0.743 at
threshold~0.1 (compared to 0.744 at association), because the location step
removes matched and proposed events in roughly equal proportion.
PyOcto+NonLinLoc likewise maintains an \ac{FDR} of 0.781, identical to the
association stage, indicating that virtually all of PyOcto's proposed event
hypotheses survive NonLinLoc inversion. This persistence suggests that the
proposed events possess internally consistent pick sets and plausible
hypocentral solutions, strengthening the hypothesis that many represent genuine
seismicity below the OGS catalog completeness threshold rather than spurious
associations. As thresholds rise to 0.3, \ac{FDR} decreases to 0.658
(GaMMA+NonLinLoc) and 0.675 (PyOcto+NonLinLoc), confirming that threshold
tuning remains the primary lever for controlling the proposed-event volume in
the final catalog.

\paragraph{Spatial residuals.}
The epicentral distance and depth difference histograms
(Table~\ref{tab:pyocto_nll1d_spatial}) reveal that PyOcto+NonLinLoc achieves 
comparable spatial accuracy to GaMMA+NonLinLoc for the events that both locate, 
but crucially does so for a much larger event set:
\begin{itemize}
  \item \textbf{Epicentral distance (0.1):} median~0.72\,km (mean~1.17\,km, 
    RMSE~1.81\,km), with the 90th percentile at 2.75\,km. These values are 
    slightly larger than GaMMA's (median~0.67\,km) because PyOcto locates 86 
    additional events, many of which are smaller or more peripheral events with 
    inherently larger location uncertainties.
  \item \textbf{Depth (0.1):} mean~$0.15$\,km (median~$0.68$\,km,
    std~3.22\,km). The interquartile range ($-0.86$\,km to $+1.89$\,km) is 
    wider than GaMMA's, again reflecting the inclusion of harder-to-locate 
    events.
  \item \textbf{Threshold 0.3:} Spatial accuracy improves: median epicentral 
    distance decreases to 0.62\,km (mean~0.99\,km), and depth standard 
    deviation narrows to 2.76\,km, approaching GaMMA+NonLinLoc levels.
\end{itemize}

PyOcto's velocity-model-aware association ensures that retained picks are
spatially coherent from the outset, so NonLinLoc's inversion converges for 
virtually every event hypothesis. The slightly larger spatial residuals 
compared to GaMMA are a direct consequence of PyOcto's higher event retention.
\begin{table}[H]
  \centering
  \begin{tabular}{c|c|c|}
    & Confusion Matrix & Time Residual Histogram \\\hline
    0.1 & \includegraphics[width=0.35\textwidth]{OGSCatalog/PyOcto0.1/img/.all_OGSLocalMagnitude_EventsConfMtx.pdf} & \includegraphics[width=0.55\textwidth]{OGSCatalog/PyOcto0.1/img/.all_OGSLocalMagnitude_EventsTimeDiff.pdf} \\\hline
    0.2 & \includegraphics[width=0.35\textwidth]{OGSCatalog/PyOcto0.2/img/.all_OGSLocalMagnitude_EventsConfMtx.pdf} & \includegraphics[width=0.55\textwidth]{OGSCatalog/PyOcto0.2/img/.all_OGSLocalMagnitude_EventsTimeDiff.pdf} \\\hline
    0.3 & \includegraphics[width=0.35\textwidth]{OGSCatalog/PyOcto0.3/img/.all_OGSLocalMagnitude_EventsConfMtx.pdf} & \includegraphics[width=0.55\textwidth]{OGSCatalog/PyOcto0.3/img/.all_OGSLocalMagnitude_EventsTimeDiff.pdf} \\\hline
  \end{tabular}
  \caption{[PhaseNet (INSTANCE) | PyOcto | NonLinLoc 1D] Event detection and 
  origin time accuracy using 1D velocity model, compared against the OGS 
  reference catalog.}
  \label{tab:pyocto_nll1d_events}
\end{table}

\newpage
\begin{landscape}
  \begin{table}[H]
    \centering
    \begin{tabular}{c|c|c|}
      & Epicentral Distance Difference & Depth Difference \\\hline
      0.1 & \includegraphics[width=0.55\textwidth]{OGSCatalog/PyOcto0.1/img/.all_OGSLocalMagnitude_EpiDistDiff.pdf} & \includegraphics[width=0.55\textwidth]{OGSCatalog/PyOcto0.1/img/.all_OGSLocalMagnitude_DepthDiff.pdf} \\\hline
      0.2 & \includegraphics[width=0.55\textwidth]{OGSCatalog/PyOcto0.2/img/.all_OGSLocalMagnitude_EpiDistDiff.pdf} & \includegraphics[width=0.55\textwidth]{OGSCatalog/PyOcto0.2/img/.all_OGSLocalMagnitude_DepthDiff.pdf} \\\hline
      0.3 & \includegraphics[width=0.55\textwidth]{OGSCatalog/PyOcto0.3/img/.all_OGSLocalMagnitude_EpiDistDiff.pdf} & \includegraphics[width=0.55\textwidth]{OGSCatalog/PyOcto0.3/img/.all_OGSLocalMagnitude_DepthDiff.pdf} \\\hline
    \end{tabular}
    \caption{[PhaseNet (INSTANCE) | PyOcto | NonLinLoc 1D] Epicentral distance 
    and depth differences as a function of P and S thresholds.}
    \label{tab:pyocto_nll1d_spatial}
  \end{table}
\end{landscape}

\newpage
\section{Magnitude Estimation}
With the location stage revealing a clear advantage for the PyOcto pathway in
event retention, we now assess magnitude estimation by comparing the local
magnitudes computed by each PhaseNet (INSTANCE)+Associator+NonLinLoc
combination against the OGS reference catalog. For each matched event the 
magnitude residual is
\[
  \Delta M_L = M_{L,\,\text{detected}} - M_{L,\,\text{catalog}}.
\]

Because magnitude estimation acts on located events, its reach is directly 
governed by the NonLinLoc recall. For GaMMA+NonLinLoc (280~$MH_E$ at threshold 
0.1), magnitudes are computed on a substantially smaller subset of the catalog 
than for PyOcto+NonLinLoc (366~$MH_E$).

\subsection{GaMMA}
The left column of Table~\ref{tab:gamma_magnitude} shows the magnitude
distribution of matched events compared to the OGS catalog, while the right
column displays the magnitude residuals. At threshold 0.1, the 280 matched 
events yield an overall magnitude residual of mean~$\Delta M_L = -0.193$ (MAE 
$= 0.225$, RMSE $= 0.316$, median~$= -0.108$). The negative mean indicates a 
systematic underestimation of magnitudes by the automated pipeline, most
pronounced below $M_L \approx 2$.

The residual scatter plots in the right column reveal a
magnitude-dependent pattern:
\begin{itemize}
  \item For $M_L \gtrsim 2$, automated and catalog magnitudes agree well, with 
    residuals concentrated within $\pm 0.3$ magnitude units.
  \item Below $M_L \approx 2$, the residual distribution broadens considerably 
    and exhibits a bimodal structure, with one cluster near zero and another 
    showing systematic underestimation by 0.3--0.8 magnitude units.
\end{itemize}
This bimodality likely arises from differences in the station subsets used by 
the automated and manual analyses.

At thresholds 0.2 and 0.3, the magnitude residual statistics remain stable 
(mean $\Delta M_L \approx -0.18$ to $-0.19$, MAE $\approx 0.22$).
\begin{table}[H]
  \centering
  \begin{tabular}{c|c|c|}
    & Magnitude distribution & Magnitude residuals \\\hline
    0.1 & \includegraphics[width=0.3\textwidth]{OGSCatalog/GaMMA0.1/img/.all_OGSLocalMagnitude_MagLDist.pdf} & \includegraphics[width=0.6\textwidth]{OGSCatalog/GaMMA0.1/img/.all_OGSLocalMagnitude_MagLDiff.pdf} \\\hline
    0.2 & \includegraphics[width=0.3\textwidth]{OGSCatalog/GaMMA0.2/img/.all_OGSLocalMagnitude_MagLDist.pdf} & \includegraphics[width=0.6\textwidth]{OGSCatalog/GaMMA0.2/img/.all_OGSLocalMagnitude_MagLDiff.pdf} \\\hline
    0.3 & \includegraphics[width=0.3\textwidth]{OGSCatalog/GaMMA0.3/img/.all_OGSLocalMagnitude_MagLDist.pdf} & \includegraphics[width=0.6\textwidth]{OGSCatalog/GaMMA0.3/img/.all_OGSLocalMagnitude_MagLDiff.pdf} \\\hline
  \end{tabular}
  \caption{[PhaseNet (INSTANCE) | GaMMA | NonLinLoc 1D] Magnitude estimation 
  performance across detection thresholds.}
  \label{tab:gamma_magnitude}
\end{table}

\subsection{PyOcto}
Table~\ref{tab:pyocto_magnitude} presents the magnitude estimation results for 
the PyOcto+NonLinLoc pathway, which benefits from a substantially larger 
matched-event set. At threshold 0.1, the 366~$MH_E$ yield 
mean~$\Delta M_L = -0.190$ (MAE $= 0.224$, RMSE $= 0.313$, median~$= -0.117$).

The magnitude distribution plots in the left column show that PyOcto recovers a 
more complete magnitude--frequency distribution, including more events in the 
$M_L < 1$ range where GaMMA's location failures create gaps. The residual 
scatter plots in the right column display the same magnitude-dependent pattern 
as GaMMA: tight agreement above $M_L \approx 2$ and increasing scatter below, 
with the characteristic bimodal structure at low magnitudes.

At higher thresholds, the PyOcto pathway shows slight improvement in residual 
statistics: mean~$\Delta M_L$ improves from $-0.190$ (threshold 0.1) to 
$-0.174$ (threshold 0.3), with MAE decreasing from $0.224$ to $0.218$ and 
median from $-0.117$ to $-0.099$. This marginal improvement reflects the 
preferential removal of poorly constrained low-magnitude events at higher 
thresholds, which have the largest magnitude residuals.
\begin{table}[H]
  \centering
  \begin{tabular}{c|c|c|}
    & Magnitude distribution & Magnitude residuals \\\hline
    0.1 & \includegraphics[width=0.3\textwidth]{OGSCatalog/PyOcto0.1/img/.all_OGSLocalMagnitude_MagLDist.pdf} & \includegraphics[width=0.6\textwidth]{OGSCatalog/PyOcto0.1/img/.all_OGSLocalMagnitude_MagLDiff.pdf} \\\hline
    0.2 & \includegraphics[width=0.3\textwidth]{OGSCatalog/PyOcto0.2/img/.all_OGSLocalMagnitude_MagLDist.pdf} & \includegraphics[width=0.6\textwidth]{OGSCatalog/PyOcto0.2/img/.all_OGSLocalMagnitude_MagLDiff.pdf} \\\hline
    0.3 & \includegraphics[width=0.3\textwidth]{OGSCatalog/PyOcto0.3/img/.all_OGSLocalMagnitude_MagLDist.pdf} & \includegraphics[width=0.6\textwidth]{OGSCatalog/PyOcto0.3/img/.all_OGSLocalMagnitude_MagLDiff.pdf} \\\hline
  \end{tabular}
  \caption{[PhaseNet (INSTANCE) | PyOcto | NonLinLoc 1D] Magnitude estimation 
  performance across detection thresholds.}
  \label{tab:pyocto_magnitude}
\end{table}

\newpage
\section{Missing Picks and Events}
\subsection{Magnitude Distribution}
To further analyze which events the pipeline fails to detect, we examine the
magnitude distributions of missed (MS) and proposed (PS) events 
(Table~\ref{tab:msps_magnitude}). A systematic understanding of which events
are lost, and where new detections appear, is essential for assessing the
operational completeness of the automated system. For the matching, we use a 
tolerance of 2 seconds in origin time and 8\,km in epicentral distance.

\paragraph{GaMMA+NonLinLoc missed events.}
At threshold 0.1, GaMMA+NonLinLoc misses 100 catalog events, distributed as
follows: 46 with $M_L < 1$, 43 with $1 \leq M_L < 2$, 11 with $2 \leq M_L < 3$.
At threshold 0.3, the missed count is 102, with 50 below $M_L = 1$, 45 between 
$M_L = 1$ and $M_L = 2$, and 7 with $M_L \geq 2$. The near-constancy of the 
missed-event count across thresholds suggests that GaMMA's location failures 
are the dominant source of loss, not the picker threshold.

\paragraph{PyOcto+NonLinLoc missed events.}
PyOcto misses far fewer events: only 14 at threshold 0.1 (6 with $M_L < 1$, 2 
with $1 \leq M_L < 2$, 6 with $2 \leq M_L < 3$). At threshold 0.3, $MS_E$ rise 
to 27 (11 with $M_L < 1$, 8 with $1 \leq M_L < 2$, 8 with $2 \leq M_L < 3$). A 
verification of all events with $M_L > 2$ confirms their inclusion in the ML 
catalog; however, they occur at distances exceeding the threshold defined for 
the matching criteria and lie near the boundary of the study area.

\paragraph{Proposed events.}
Proposed events (1,304 for PyOcto; 810 for GaMMA, at threshold 0.1) are 
distributed throughout the monitoring area, with concentrations that may 
indicate genuine micro-seismicity clusters below the OGS catalog detection 
threshold. The spatial distributions of these missed and proposed events are 
shown in Figs.~\ref{fig:gamma_eventsFalse} and~\ref{fig:pyocto_eventsFalse}. A 
spurious event with $M_L \geq 10$ located near the center of the network by 
PyOcto+NonLinLoc at threshold 0.1 is a clear outlier in the magnitude 
distribution and is likely a misassociation or a processing artifact, as no 
such event exists in the GaMMA+NonLinLoc results. This event should be flagged 
for manual review.

\begin{landscape}
  \begin{table}[H]
    \centering
    \begin{tabular}{c|c|c|}
      & [GaMMA | NonLinLoc 1D] & [PyOcto | NonLinLoc 1D] \\\hline
      0.1 & \includegraphics[width=0.5\textwidth]{OGSCatalog/GaMMA0.1/img/.all_OGSLocalMagnitude_MSPSMagLDist.pdf} & \includegraphics[width=0.5\textwidth]{OGSCatalog/PyOcto0.1/img/.all_OGSLocalMagnitude_MSPSMagLDist.pdf} \\\hline
      0.2 & \includegraphics[width=0.5\textwidth]{OGSCatalog/GaMMA0.2/img/.all_OGSLocalMagnitude_MSPSMagLDist.pdf} & \includegraphics[width=0.5\textwidth]{OGSCatalog/PyOcto0.2/img/.all_OGSLocalMagnitude_MSPSMagLDist.pdf} \\\hline
      0.3 & \includegraphics[width=0.5\textwidth]{OGSCatalog/GaMMA0.3/img/.all_OGSLocalMagnitude_MSPSMagLDist.pdf} & \includegraphics[width=0.5\textwidth]{OGSCatalog/PyOcto0.3/img/.all_OGSLocalMagnitude_MSPSMagLDist.pdf} \\\hline
    \end{tabular}
    \caption{Magnitude distribution of missed (MS) and proposed (PS) events 
    as a function of P and S thresholds. MS events are OGS catalog entries not 
    matched by the automated system; PS events are automated detections absent 
    from the OGS catalog.}
    \label{tab:msps_magnitude}
  \end{table}
\end{landscape}

\newpage
\begin{figure}
  \centering
  \includegraphics[width=.9\textwidth]{OGSCatalog/GaMMA0.1/img/.all_OGSLocalMagnitude_EventsFalse.pdf}
  \caption{[PhaseNet (INSTANCE) | GaMMA | NonLinLoc 1D] Spatial distribution 
  of missed (blue) and proposed (pink) events at detection threshold 0.1.}
  \label{fig:gamma_eventsFalse}
\end{figure}

\begin{figure}
  \centering
  \includegraphics[width=.9\textwidth]{OGSCatalog/PyOcto0.1/img/.all_OGSLocalMagnitude_EventsFalse.pdf}
  \caption{[PhaseNet (INSTANCE) | PyOcto | NonLinLoc 1D] Spatial distribution 
  of missed (blue) and proposed (pink) events at detection threshold 0.1.
  PyOcto misses only 14 events concentrated at the periphery of the Bullettin 
  area (red polygon). Proposed events (1,304) are distributed throughout the 
  monitoring area, with concentrations that may indicate genuine 
  micro-seismicity clusters below the OGS catalog detection threshold.}
  \label{fig:pyocto_eventsFalse}
\end{figure}

\begin{comment}
  
\newpage
\section{Risk Management Implications}
From a risk-management perspective, the choice of picker thresholds and
association parameters determines the balance between catalog recall
(lower MS) and operational FDR (lower PS). The analysis in this
chapter suggests two practical operating modes:
\begin{itemize}
  \item Surveillance mode (24/7 monitoring): moderate thresholds (e.g., P/S
    in the 0.2--0.3 range) with PyOcto or GaMMA tuned for low PS. This
    minimizes false alarms and analyst load while maintaining sensitivity
    to felt events and most M\,$\gtrsim$\,2 microseismicity.
  \item Response mode (after felt mainshock or in swarm): relaxed thresholds
    to prioritize recall and situational awareness, with automated
    QC (e.g., minimum station count, timing residual gates) to filter the
    highest-confidence detections for rapid dissemination.
\end{itemize}

Timing precision (narrow residuals) directly improves phase association and
location, leading to more reliable impact assessment (e.g., faster, more
accurate epicentral mapping to support civil protection decisions). Magnitude
bias and scatter influence downstream hazard indicators (frequency--magnitude
statistics and aftershock productivity), so periodic calibration against the
OGS catalog is recommended.

\subsection{Operational Recommendations}
\begin{itemize}
  \item Picker: Use PhaseNet trained on INSTANCE as default; validate new
    deployments with time-residual diagnostics and MS/PS distributions.
  \item Associator: Prefer PyOcto in dense local networks or swarm scenarios;
    prefer GaMMA where probabilistic treatment of uncertainty is
    advantageous (sparser coverage, noisier picks).
  \item Thresholds: Start at 0.2--0.3 for P and S; adjust per-station or
    per-region if systematic MS/PS patterns appear in quality-control
    dashboards.
  \item QC gates: Enforce minimum station count, maximum per-station residual,
    and P/S ordering checks before publishing events.
  \item Review loop: Track KPIs (recall, FDR, lead time) and re-tune monthly
    or after major configuration changes.
\end{itemize}

\subsection{Limitations and Future Work}
Integration of other associators (e.g., REAL) and consistent 3D velocity models 
across all workflows are active areas of development. Magnitude calibration 
against multiple reference scales and robust uncertainty quantification (e.g., 
ensemble pickers/associators) will further enhance reliability for operational 
risk management.

\paragraph{Comparative summary.} For day-to-day surveillance, PhaseNet (INSTANCE)
plus PyOcto or GaMMA at P/S thresholds of 0.2--0.3 provides a balanced trade-off
between recall and FDR, with PyOcto often preferred in dense
local networks and GaMMA in noisier or sparser conditions. Incorporating 3D
velocity models generally improves event matching and timing residuals across
all workflows.

\begin{table}[H]
  \centering
  \caption{At-a-glance qualitative comparison of association/location workflows 
  at typical P/S thresholds (0.2--0.3). Categories summarize observed behavior 
  in this study; use KPI dashboards to refine per-network settings.}
  \label{tab:algo_comparison}
  \begin{tabular}{l|c|c|c|l}
    Algorithm & Latency & Recall & FDR & Best for \\\hline
    GaMMA  & Medium & High & Low--Medium & Sparse/noisy networks; probabilistic uncertainty handling \\
    PyOcto & Low    & Medium--High & Low & Dense local networks; swarm sequences \\
  \end{tabular}
\end{table}
\end{comment}
